% ===========================================================================
% MANUAL TÉCNICO - TEST VOCACIONAL
% Proyecto: test-vocacional
% Autor: [Equipo de Desarrollo]
% Fecha: Abril 2025
% ===========================================================================

\documentclass[12pt, a4paper]{report}

% ---------------------------------------------------------------------------
% PAQUETES BÁSICOS
% ---------------------------------------------------------------------------
\usepackage[utf8]{inputenc}
\usepackage[T1]{fontenc}
\usepackage[spanish]{babel}
\usepackage{geometry}
\usepackage{graphicx}
\usepackage{hyperref}
\usepackage{listings}
\usepackage{xcolor}
\usepackage{fancyhdr}
\usepackage{titlesec}
\usepackage{tocloft}
\usepackage{tabularx}
\usepackage{longtable}
\usepackage{multirow}
\usepackage{booktabs}
\usepackage{amsmath}
\usepackage{amssymb}
\usepackage{enumitem}
\usepackage{verbatim}
\usepackage{float}
\usepackage{tikz}
\usepackage{pgfplots}
\usepackage{minted}
\usepackage{caption}
\usepackage{subcaption}

% ---------------------------------------------------------------------------
% CONFIGURACIÓN DE PÁGINA
% ---------------------------------------------------------------------------
\geometry{
    top=3cm,
    bottom=2.5cm,
    left=2.5cm,
    right=2.5cm
}

\pagestyle{fancy}
\fancyhf{}
\fancyhead[L]{\leftmark}
\fancyhead[R]{\thepage}
\fancyfoot[C]{}
\renewcommand{\headrulewidth}{0.5pt}

% ---------------------------------------------------------------------------
% COLORES PARA CÓDIGO
% ---------------------------------------------------------------------------
\definecolor{codegreen}{rgb}{0,0.6,0}
\definecolor{codegray}{rgb}{0.5,0.5,0.5}
\definecolor{codepurple}{rgb}{0.58,0,0.82}
\definecolor{backcolour}{rgb}{0.95,0.95,0.92}

% ---------------------------------------------------------------------------
% LISTINGS CONFIGURACIÓN
% ---------------------------------------------------------------------------
\lstdefinestyle{mystyle}{
    backgroundcolor=\color{backcolour},
    commentstyle=\color{codegreen},
    keywordstyle=\color{magenta},
    numberstyle=\tiny\color{codegray},
    stringstyle=\color{codepurple},
    basicstyle=\ttfamily\footnotesize,
    breakatwhitespace=false,
    breaklines=true,
    captionpos=b,
    keepspaces=true,
    numbers=left,
    numbersep=5pt,
    showspaces=false,
    showstringspaces=false,
    showtabs=false,
    tabsize=2,
    frame=single,
    framerule=0.6pt
}

\lstset{style=mystyle}

% Configuración adicional para listings de estructura de directorios
\lstdefinestyle{dirstyle}{
    backgroundcolor=\color{backcolour},
    basicstyle=\ttfamily\footnotesize,
    breaklines=true,
    breakatwhitespace=true,
    frame=single,
    framerule=0.6pt,
    numbers=none,
    showstringspaces=false,
    language=bash
}

% ---------------------------------------------------------------------------
% CONFIGURACIÓN DE SECCIONES
% ---------------------------------------------------------------------------
\titleformat{\chapter}[display]
{\normalfont\huge\bfseries}{\chaptertitlename\ \thechapter}{20pt}{\Huge}
\titleformat{\section}
{\normalfont\Large\bfseries}{\thesection}{1em}{}
\titleformat{\subsection}
{\normalfont\large\bfseries}{\thesubsection}{1em}{}

% ---------------------------------------------------------------------------
% META DATOS DEL PDF
% ---------------------------------------------------------------------------
\hypersetup{
    colorlinks=true,
    linkcolor=blue,
    filecolor=magenta,
    urlcolor=cyan,
    pdftitle={Manual Técnico - Test Vocacional},
    pdfauthor={Equipo de Desarrollo},
    pdfsubject={Documentación Técnica},
    pdfkeywords={TestVocacional, RIASEC, PHP, Manual Técnico}
}

% ---------------------------------------------------------------------------
% ESTILOS PARA DIAGRAMAS TIKZ (Corrección de estilos no definidos)
% ---------------------------------------------------------------------------
\tikzset{
    entity/.style={rectangle, draw, rounded corners, fill=blue!10, minimum width=2cm, minimum height=1cm, align=center},
    relation/.style={->, thick, >=stealth},
    attribute/.style={ellipse, draw, fill=yellow!10}
}

% ===========================================================================
% INICIO DEL DOCUMENTO
% ===========================================================================

\begin{document}

% ---------------------------------------------------------------------------
% PORTADA
% ---------------------------------------------------------------------------
\begin{titlepage}
    \centering
    \vspace*{2cm}
    
    {\Huge \textbf{MANUAL TÉCNICO}}\\[0.5cm]
    {\Huge \textbf{TEST VOCACIONAL}}\\[2cm]
    
    % Imagen condicional para evitar error si no existe
   
    
    {\Large \textbf{Proyecto: test-vocacional}}\\[0.5cm]
    {\large Arquitectura Full-Stack MVC}\\[1.5cm]
    
    \rule{\textwidth}{0.4pt}\\[1cm]
    {\large \textbf{Stack Tecnológico}}\\
    {\large PHP 8.2+ · MySQL · HTML5/CSS3 · JavaScript}\\[1cm]
    {\large \textbf{Modelo RIASEC (Holland)}}\\[1cm]
    \rule{\textwidth}{0.4pt}\\[2cm]
    
    {\large Versión 1.0}\\[0.5cm]
    {\large Abril 2025}\\[3cm]
    
    {\normalsize \textbf{Este documento contiene:}}\\
    {\normalsize Documentación técnica completa de arquitectura,}\\
    {\normalsize configuración, despliegue y mantenimiento.}
    
    \vfill
\end{titlepage}

% ---------------------------------------------------------------------------
% PÁGINA DE AUTORES
% ---------------------------------------------------------------------------
\newpage
\thispagestyle{empty}
\vspace*{5cm}
\begin{flushleft}
    {\large \textbf{Equipo de Desarrollo}}\\[0.5cm]
    Rubén Jaramillo / Miguel Ulcuango \\[0.2cm]
    Contacto: info@desarrollodesoftware.com.ec \\[2cm]
    
    {\large \textbf{Ubicación del Proyecto}}\\[0.5cm]
    \texttt{/Users/test-vocacional} \\[2cm]
\end{flushleft}

% ---------------------------------------------------------------------------
% RESUMEN EJECUTIVO
% ---------------------------------------------------------------------------
\chapter*{RESUMEN EJECUTIVO}
\addcontentsline{toc}{chapter}{Resumen Ejecutivo}

El \textbf{Test Vocacional} es una aplicación web full-stack diseñada para evaluar las aptitudes vocacionales de estudiantes mediante el modelo RIASEC de Holland. El sistema evalúa seis categorías principales: Realista, Inventivo/Investigador, Artístico, Social, Emprendedor y Convencional.

\begin{itemize}[leftmargin=2cm]
    \item \textbf{Tipo:} Aplicación Web Full-Stack (PHP + MySQL + Frontend)
    \item \textbf{Arquitectura:} Patrón MVC personalizado
    \item \textbf{Base de Datos:} MySQL/MariaDB (5 tablas principales)
    \item \textbf{Usuarios:} Múltiples roles (administrador, zonal, DECE, estudiante)
    \item \textbf{Reportes:} Generación de PDF, Excel y documentos Word
    \item \textbf{Stack:} PHP 7.4+, HTML5, CSS3, JavaScript ES6
\end{itemize}

El sistema está diseñado para instituciones educativas y permite:
\begin{itemize}
    \item Administración de usuarios e instituciones
    \item Gestión de preguntas del test (importación vía Excel)
    \item Aplicación del test vocacional en línea
    \item Cálculo automático de perfiles RIASEC
    \item Generación de informes individuales y grupales
    \item Paneles específicos por rol de usuario
\end{itemize}

% ---------------------------------------------------------------------------
% ÍNDICE
% ---------------------------------------------------------------------------
\tableofcontents
\newpage

\listoffigures
\newpage

\listoftables
\newpage

% ===========================================================================
% CAPÍTULO 1: ARQUITECTURA GENERAL
% ===========================================================================
\chapter{ARQUITECTURA GENERAL DEL SISTEMA}

\section{Visión General}
El proyecto \texttt{test-vocacional} es una aplicación web completa que implementa un sistema de orientación vocacional basado en la teoría tipológica de John L. Holland (RIASEC). La arquitectura sigue un patrón \textbf{MVC (Model-View-Controller)} personalizado sin framework externo, lo que proporciona total control sobre la estructura y permite un fácil mantenimiento.

\section{Tecnologías Principales}

\subsection{Backend}
\begin{table}[H]
\centering
\begin{tabularx}{\textwidth}{|l|X|l|}
\hline
\textbf{Tecnología} & \textbf{Descripción} & \textbf{Versión} \\
\hline
PHP & Lenguaje de programación server-side & 7.4+ \\
MySQL & Sistema de gestión de base de datos & 5.7+/8.0 \\
PDO & Capa de abstracción de base de datos & Native \\
Composer & Gestor de dependencias PHP & 2.x \\
\hline
\end{tabularx}
\caption{Tecnologías Backend}
\end{table}

\subsection{Frontend}
\begin{table}[H]
\centering
\begin{tabularx}{\textwidth}{|l|X|l|}
\hline
\textbf{Tecnología} & \textbf{Descripción} & \textbf{Versión} \\
\hline
HTML5 & Estructura y semántica & Standard \\
CSS3 & Estilos y diseño & Standard \\
JavaScript & Lógica del lado del cliente & ES6+ \\
TCPDF & Generación de PDFs en servidor & 6.4 \\
PHPSpreadsheet & Generación de Excel & 5.1 \\
PHPWord & Generación de documentos Word & 1.4 \\
\hline
\end{tabularx}
\caption{Tecnologías Frontend}
\end{table}

\section{Patrones Arquitectónicos}

\subsection{MVC (Model-View-Controller)}
El sistema implementa una variante del patrón MVC:
\begin{itemize}
    \item \textbf{Models}: Clases que encapsulan acceso a datos (extienden \texttt{BaseModel})
    \item \textbf{Views}: Vistas PHP que renderizan HTML embebido
    \item \textbf{Controllers}: Clases que procesan requests y devuelven respuestas
\end{itemize}

\subsection{Singleton Database}
Clase \texttt{Database} implementa conexión única PDO:
\begin{lstlisting}[language=PHP, caption=Patrón Singleton Database]
class Database {
    private static $instance = null;
    private $connection;
    
    private function __construct() {
        $this->connection = new PDO(...);
    }
    
    public static function getInstance() {
        if (self::$instance == null) {
            self::$instance = new Database();
        }
        return self::$instance;
    }
}
\end{lstlisting}

\subsection{Middleware de Autenticación}
Filtro \texttt{AuthMiddleware.php} para control de acceso basado en roles.

% ===========================================================================
% CAPÍTULO 2: ESTRUCTURA DE DIRECTORIOS
% ===========================================================================
\chapter{ESTRUCTURA DE DIRECTORIOS}

\section{Organización del Proyecto}
\begin{lstlisting}[style=dirstyle, caption=Estructura de directorios del proyecto]
test-vocacional2/
+-- config/              # Configuraciones globales
|   +-- config.php      # Constantes, BD, SMTP, RIASEC
|   +-- database.php    # Clase singleton Database (PDO)
+-- controllers/        # Controladores MVC
|   +-- AdminController.php
|   +-- AuthController.php
|   +-- DECEController.php
|   +-- QuestionController.php
|   +-- TestController.php
|   +-- ZonaController.php
+-- models/             # Modelos MVC
|   +-- BaseModel.php
|   +-- User.php
|   +-- Question.php
|   +-- TestResult.php
|   +-- VocationalTest.php
|   +-- Institucion.php
+-- views/              # Vistas (páginas PHP)
|   +-- layout/
|   |   +-- header.php
|   |   +-- footer.php
|   |   +-- sidebar.php
|   +-- login.php
|   +-- register.php
|   +-- test_form.php
|   +-- test_results.php
|   +-- admin_users.php
|   +-- admin_institutions.php
|   +-- admin_questions.php
|   +-- admin_reports.php
|   +-- admin_dashboard.php
|   +-- DECE_dashboard.php
|   +-- zonal_dashboard.php
|   +-- profile.php
+-- services/           # Servicios de negocio
|   +-- UserService.php
|   +-- InstitutionService.php
|   +-- ReportService.php
+-- middleware/         # Middleware
|   +-- AuthMiddleware.php
+-- utils/              # Utilidades
|   +-- PDFGenerator.php
|   +-- ExcelGenerator.php
|   +-- EmailSender.php
|   +-- Recommendations.php
|   +-- QuestionImporter.php
|   +-- riasec_helpers.php
+-- assets/             # Recursos estáticos
|   +-- css/
|   |   +-- styles.css (~30KB)
|   |   +-- vendor/
|   +-- js/
|   |   +-- main.js
|   +-- img/
|       +-- logo.png
+-- storage/            # Logs y archivos temporales
+-- vendor/             # Dependencias Composer
+-- index.php           # Router principal
+-- .htaccess           # RewriteEngine (Apache)
+-- bdd.sql             # Base de datos completa
+-- composer.json       # Dependencias PHP
+-- package.json        # Dependencias Node.js
\end{lstlisting}

\section{Propósito de Directorios}

\begin{table}[H]
\centering
\begin{tabularx}{\textwidth}{|l|X|}
\hline
\textbf{Directorio} & \textbf{Propósito} \\
\hline
\texttt{config/} & Almacena configuraciones globales, constantes del sistema y conexión a base de datos \\
\hline
\texttt{controllers/} & Contiene los controladores que manejan las peticiones HTTP y la lógica de los endpoints \\
\hline
\texttt{models/} & Modelos que representan entidades de datos y operaciones CRUD \\
\hline
\texttt{views/} & Plantillas HTML/PHP para renderizar la interfaz de usuario \\
\hline
\texttt{services/} & Capa de servicios con lógica de negocio reutilizable \\
\hline
\texttt{middleware/} & Filtros para autenticación y autorización de rutas \\
\hline
\texttt{utils/} & Utilidades y helpers para generación de reportes y procesamiento \\
\hline
\texttt{assets/} & Archivos estáticos (CSS, JS, imágenes) para el frontend \\
\hline
\texttt{storage/} & Archivos temporales y logs del sistema \\
\hline
\texttt{vendor/} & Dependencias externas instaladas con Composer \\
\hline
\end{tabularx}
\caption{Propósito de Directorios}
\end{table}

% ===========================================================================
% CAPÍTULO 3: CONFIGURACIÓN
% ===========================================================================
\chapter{CONFIGURACIÓN DEL SISTEMA}

\section{Archivo de Configuración Principal}
\subsection{config/config.php}
\begin{lstlisting}[language=PHP, caption=config/config.php]
<?php
// Configuración de la base de datos
define('DB_HOST', 'localhost');
define('DB_NAME', 'test_vocacional');
define('DB_USER', 'root');
define('DB_PASS', '');
define('DB_CHARSET', 'utf8mb4');

// Configuración de la aplicación
define('APP_URL', 'http://localhost/test-vocacional');
define('APP_EMAIL', 'noreply@testvocacional.edu.ec');
define('SESSION_LIFETIME', 3600); // 1 hora
define('PASSWORD_MIN_LENGTH', 6);

// Configuración del test vocacional RIASEC
define('TEST_QUESTIONS_PER_CATEGORY', 10);
define('TOTAL_TEST_QUESTIONS', 60);
define('TEST_RETAKE_MONTHS', 6);

// Umbrales de puntuación RIASEC
define('APTO_THRESHOLD', 80);       // Puntaje mínimo para apto
define('POTENCIAL_THRESHOLD', 60);  // Puntaje mínimo para potencial

// Zonas geográficas (Ecuador)
define('ZONA_COSTAL', 'costal');
define('ZONA_SIERRA', 'sierra');
define('ZONA_AMAZONICA', 'amazonica');
define('ZONA_INSULAR', 'insular');

// Roles del sistema
define('ROL_ADMIN', 'administrador');
define('ROL_ZONAL', 'zonal');
define('ROL_DECE', 'dece');
define('ROL_ESTUDIANTE', 'estudiante');

// Configuración SMTP (opcional para emails)
define('SMTP_HOST', 'smtp.gmail.com');
define('SMTP_PORT', 587);
define('SMTP_USER', '');
define('SMTP_PASS', '');
define('SMTP_FROM', 'Test Vocacional <noreply@testvocacional.edu.ec>');

// Timezone
date_default_timezone_set('America/Guayaquil');
?>
\end{lstlisting}

\section{Conexión a Base de Datos}
\subsection{config/database.php}
\begin{lstlisting}[language=PHP, caption=config/database.php]
<?php
class Database {
    private static $instance = null;
    private $connection;
    
    private function __construct() {
        try {
            $dsn = "mysql:host=" . DB_HOST . ";dbname=" . DB_NAME . ";charset=" . DB_CHARSET;
            $options = [
                PDO::ATTR_ERRMODE => PDO::ERRMODE_EXCEPTION,
                PDO::ATTR_DEFAULT_FETCH_MODE => PDO::FETCH_ASSOC,
                PDO::ATTR_EMULATE_PREPARES => false,
            ];
            $this->connection = new PDO($dsn, DB_USER, DB_PASS, $options);
        } catch (PDOException $e) {
            die("Error de conexión: " . $e->getMessage());
        }
    }
    
    public static function getInstance() {
        if (self::$instance == null) {
            self::$instance = new Database();
        }
        return self::$instance;
    }
    
    public function getConnection() {
        return $this->connection;
    }
    
    public function prepare($sql) {
        return $this->connection->prepare($sql);
    }
}
?>
\end{lstlisting}

% ===========================================================================
% CAPÍTULO 4: MODELO DE DATOS
% ===========================================================================
\chapter{MODELO DE DATOS}

\section{Diagrama Entidad-Relación}
\begin{figure}[H]
\centering
\begin{tikzpicture}[
    scale=0.9,
    every node/.style={transform shape},
    entity/.style={
        rectangle, 
        draw, 
        rounded corners, 
        fill=blue!10, 
        minimum width=2.2cm, 
        minimum height=1.2cm, 
        align=center,
        font=\small
    },
    relation/.style={->, thick}
]
    % --- NODOS con coordenadas absolutas ---
    \node[entity] (usuarios) at (0,0) {usuarios};
    \node[entity] (instituciones) at (5,0) {instituciones\\educativas};
    \node[entity] (preguntas) at (0,-3) {preguntas};
    \node[entity] (detalle) at (5,-3) {respuestas\\detalle};
    \node[entity] (resultados) at (0,-6) {resultados\\test};
    
    % --- RELACIONES ---
    \draw[relation] (usuarios) -- node[above] {\small 1:N} (detalle);
    \draw[relation] (preguntas) -- node[above] {\small 1:N} (detalle);
    \draw[relation] (usuarios) -- node[right] {\small 1:N} (resultados);
    \draw[relation] (detalle) -- node[right] {\small N:1} (resultados);
    \draw[relation] (usuarios) -- node[above] {\small N:1} (instituciones);
\end{tikzpicture}
\caption{Diagrama ER simplificado del sistema}
\label{fig:diagrama-er}
\end{figure}

\section{Tablas Principales}

\subsection{Tabla: usuarios}
\begin{table}[H]
\centering
\begin{tabularx}{\textwidth}{|l|X|l|}
\hline
\textbf{Campo} & \textbf{Tipo} & \textbf{Descripción} \\
\hline
id & INT(11) PK AI & Identificador único \\
cedula & VARCHAR(20) UNIQUE & Número de cédula \\
nombre & VARCHAR(100) & Nombre completo \\
email & VARCHAR(100) UNIQUE & Correo electrónico \\
password & VARCHAR(255) & Contraseña hasheada \\
rol & ENUM('administrador','zonal','dece','estudiante') & Rol del usuario \\
institucion\_id & INT(11) FK & Referencia a institución \\
zona & VARCHAR(50) & Zona geográfica \\
activo & TINYINT(1) & Estado (activo/inactivo) \\
created\_at & TIMESTAMP & Fecha de creación \\
updated\_at & TIMESTAMP & Fecha de actualización \\
\hline
\end{tabularx}
\caption{Estructura de la tabla usuarios}
\end{table}

\subsection{Tabla: preguntas}
\begin{table}[H]
\centering
\begin{tabularx}{\textwidth}{|l|X|l|}
\hline
\textbf{Campo} & \textbf{Tipo} & \textbf{Descripción} \\
\hline
id & INT(11) PK AI & Identificador único \\
texto & TEXT & Texto de la pregunta \\
categoria & VARCHAR(50) & Categoría RIASEC \\
tipo & ENUM('intereses','habilidades','valores') & Tipo de pregunta \\
orden & INT(11) & Orden de aparición \\
activo & TINYINT(1) & Estado \\
created\_at & TIMESTAMP & Fecha creación \\
\hline
\end{tabularx}
\caption{Estructura de la tabla preguntas}
\end{table}

\subsection{Tabla: respuestas\_detalle}
\begin{table}[H]
\centering
\begin{tabularx}{\textwidth}{|l|X|l|}
\hline
\textbf{Campo} & \textbf{Tipo} & \textbf{Descripción} \\
\hline
id & INT(11) PK AI & Identificador único \\
test\_id & INT(11) FK & Referencia al test \\
pregunta\_id & INT(11) FK & Referencia a pregunta \\
respuesta & INT(1) & Valor 1-5 (Likert) \\
\hline
\end{tabularx}
\caption{Estructura de respuestas detalle}
\end{table}

\subsection{Tabla: resultados\_test}
\begin{table}[H]
\centering
\begin{tabularx}{\textwidth}{|l|X|l|}
\hline
\textbf{Campo} & \textbf{Tipo} & \textbf{Descripción} \\
\hline
id & INT(11) PK AI & Identificador \\
test\_id & INT(11) FK & Referencia test \\
usuario\_id & INT(11) FK & Usuario \\
resultados & JSON & Puntuaciones RIASEC en JSON \\
fecha & DATETIME & Fecha del test \\
tiempo\_empleado & INT & Tiempo en segundos \\
\hline
\end{tabularx}
\caption{Estructura de resultados de test}
\end{table}

\subsection{Tabla: instituciones\_educativas}
\begin{table}[H]
\centering
\begin{tabularx}{\textwidth}{|l|X|l|}
\hline
\textbf{Campo} & \textbf{Tipo} & \textbf{Descripción} \\
\hline
id & INT(11) PK AI & Identificador \\
nombre & VARCHAR(255) & Nombre institución \\
codigo & VARCHAR(50) UNIQUE & Código único \\
tipo & VARCHAR(50) & Tipo (pública/privada) \\
direccion & TEXT & Dirección \\
telefono & VARCHAR(20) & Teléfono \\
email & VARCHAR(100) & Correo \\
representante & VARCHAR(100) & Representante \\
zona & VARCHAR(50) & Zona geográfica \\
activo & TINYINT(1) & Estado \\
created\_at & TIMESTAMP & Fecha creación \\
\hline
\end{tabularx}
\caption{Estructura de instituciones}
\end{table}

% ===========================================================================
% CAPÍTULO 5: MODELO RIASEC
% ===========================================================================
\chapter{MODELO RIASEC (Holland)}

\section{Fundamento Teórico}
El modelo RIASEC de John L. Holland clasifica las vocaciones en seis categorías:
\begin{enumerate}[leftmargin=2cm]
    \item \textbf{R - Realista}: Actividades prácticas, físicas, manuales (ej: ingeniería, construcción)
    \item \textbf{I - Investigador}: Actividades intelectuales, analíticas, científicas
    \item \textbf{A - Artístico}: Actividades creativas, expresivas, estéticas
    \item \textbf{S - Social}: Actividades de servicio, enseñanza, ayuda
    \item \textbf{E - Emprendedor}: Actividades de liderazgo, persuasivas, comerciales
    \item \textbf{C - Convencional}: Actividades organizativas, administrativas, detallistas
\end{enumerate}

\section{Implementación en el Sistema}
Cada categoría evalúa tres dimensiones:
\begin{itemize}
    \item \textbf{Intereses}: Actividades que disfrutaría hacer
    \item \textbf{Habilidades}: Capacidades naturales
    \item \textbf{Valores}: Principios y creencias laborales
\end{itemize}

\section{Cálculo de Puntuaciones}
\begin{lstlisting}[language=PHP, caption=Algoritmo de cálculo RIASEC]
class VocationalTest {
    private $categories = ['R', 'I', 'A', 'S', 'E', 'C'];
    
    public function calculate($answers) {
        $scores = [];
        foreach ($this->categories as $cat) {
            $scores[$cat] = 0;
        }
        
        foreach ($answers as $answer) {
            $question = Question::find($answer['pregunta_id']);
            $category = $question->categoria; // R, I, A, S, E, C
            $value = $answer['respuesta'];     // 1-5 Likert
            $scores[$category] += $value;
        }
        
        // Normalizar a 100%
        $total = array_sum($scores);
        foreach ($scores as &$score) {
            $score = round(($score / $total) * 100, 2);
        }
        
        return $scores;
    }
}
\end{lstlisting}

\section{Interpretación de Resultados}
Según los puntajes, el sistema determina:
\begin{itemize}
    \item \textbf{Apto}: Puntaje ≥ 80\% en al menos una categoría
    \item \textbf{Potencial}: Puntaje ≥ 60\% en al menos una categoría
    \item \textbf{En desarrollo}: Puntaje < 60\%
\end{itemize}

% ===========================================================================
% CAPÍTULO 6: CONTROLADORES
% ===========================================================================
\chapter{CONTROLADORES (Controllers)}

\section{Controller Base}
Todos los controladores extienden de una clase base con métodos comunes:
\begin{itemize}
    \item \textbf{redirect()}: Redirecciones
    \item \textbf{view()}: Renderizado de vistas
    \item \textbf{json()}: Respuestas JSON
    \item \textbf{authCheck()}: Verificación de autenticación
\end{itemize}

\section{AuthController.php}
Maneja registro, login y recuperación de contraseñas.
\begin{lstlisting}[language=PHP, caption=AuthController.php - Métodos principales]
class AuthController {
    // POST /register
    public function register() {
        Validación de datos → User::createUser() → envío email → login()
    }
    
    // POST /login
    public function login() {
        Verificar credenciales → crear sesión → redirect según rol
    }
    
    // POST /logout
    public function logout() {
        Destruir sesión → redirect a login
    }
    
    // GET/POST /recover
    public function recoverPassword() {
        Generar token → enviar email de recuperación
    }
}
\end{lstlisting}

\section{TestController.php}
Gestiona el test vocacional y resultados.
\begin{lstlisting}[language=PHP, caption=TestController.php - Flujo del test]
class TestController {
    // GET /test
    public function showTest() {
        Cargar preguntas filtradas por categoría → mostrar formulario
    }
    
    // POST /test/submit
    public function submit() {
        Validar respuestas → VocationalTest::calculate()
        → Guardar resultados → Generar PDF
        → Redirect /results
    }
    
    // GET /results
    public function showResults() {
        Obtener resultados del usuario → Calcular recomendaciones
        → Mostrar gráfico RIASEC → Enlaces a ReportService
    }
}
\end{lstlisting}

\section{AdminController.php}
Panel de administración completo.
\begin{lstlisting}[language=PHP, caption=AdminController.php - Funcionalidades]
class AdminController {
    // GET /admin/dashboard
    public function dashboard() {
        Estadísticas globales: usuarios, tests, instituciones
    }
    
    // CRUD /admin/users
    public function users() {
        Listar, crear, editar, eliminar, activar/desactivar usuarios
    }
    
    // CRUD /admin/institutions
    public function institutions() {
        Gestión completa de instituciones educativas
    }
    
    // CRUD /admin/questions
    public function questions() {
        Gestión del banco de preguntas RIASEC
        Importación desde Excel → QuestionImporter
    }
    
    // Reports /admin/reports/*
    public function reports() {
        Reportes individuales y grupales via ReportService
        Exportación PDF/Excel
    }
}
\end{lstlisting}

\section{DECEController.php - ZonaController.php}
Paneles específicos con funcionalidades limitadas.

% ===========================================================================
% CAPÍTULO 7: MODELOS (Models)
% ===========================================================================
\chapter{MODELOS (Models)}

\section{BaseModel.php}
Clase base con operaciones CRUD genéricas.
\begin{lstlisting}[language=PHP, caption=BaseModel.php]
abstract class BaseModel {
    protected $db;
    protected $table;
    protected $primaryKey = 'id';
    
    public function __construct() {
        $this->db = Database::getInstance()->getConnection();
    }
    
    public function findAll($conditions = [], $orderBy = null) {
        $sql = "SELECT * FROM {$this->table}";
        // ... WHERE conditions
        // ... ORDER BY
    }
    
    public function findById($id) {
        $stmt = $this->db->prepare("SELECT * FROM {$this->table}
                                    WHERE {$this->primaryKey} = ?");
        $stmt->execute([$id]);
        return $stmt->fetch();
    }
    
    public function save($data) {
        if(isset($data[$this->primaryKey])) {
            return $this->update($data);
        } else {
            return $this->insert($data);
        }
    }
    
    public function delete($id) {
        $stmt = $this->db->prepare("DELETE FROM {$this->table}
                                    WHERE {$this->primaryKey} = ?");
        return $stmt->execute([$id]);
    }
}
\end{lstlisting}

\section{Modelos Específicos}

\subsection{User.php}
\begin{lstlisting}[language=PHP]
class User extends BaseModel {
    protected $table = 'usuarios';
    
    public function findByUsername($username) {
        $stmt = $this->db->prepare("SELECT * FROM usuarios 
                                    WHERE email = ? OR cedula = ?");
        $stmt->execute([$username, $username]);
        return $stmt->fetch();
    }
    
    public function verifyPassword($password, $hash) {
        return password_verify($password, $hash);
    }
    
    public function hashPassword($password) {
        return password_hash($password, PASSWORD_BCRYPT);
    }
    
    public function getByRole($rol) {
        return $this->findAll(['rol' => $rol]);
    }
    
    public function canTakeTest() {
        // Verifica si han pasado 6 meses desde último test
    }
}
\end{lstlisting}

\subsection{VocationalTest.php}
Manejo de tests y cálculos RIASEC.
\begin{lstlisting}[language=PHP]
class VocationalTest extends BaseModel {
    protected $table = 'resultados_test';
    
    public function calculateResults($answers) {
        $categories = ['R','I','A','S','E','C'];
        $scores = array_fill_keys($categories, 0);
        
        foreach($answers as $answer) {
            $question = Question::find($answer['pregunta_id']);
            $cat = $question['categoria'];
            $scores[$cat] += $answer['respuesta'];
        }
        
        // Normalización y clasificación
        return $this->normalizeScores($scores);
    }
    
    public function getRecommendations($scores) {
        // Lógica de recomendaciones basada en puntajes
        $maxCat = array_search(max($scores), $scores);
        return $this->generateProfile($maxCat);
    }
}
\end{lstlisting}

\subsection{Question.php}
\begin{lstlisting}[language=PHP]
class Question extends BaseModel {
    protected $table = 'preguntas';
    
    public function getByCategory($category) {
        return $this->findAll(['categoria' => $category, 'activo' => 1]);
    }
    
    public function importFromExcel($filePath) {
        // Importación masiva de preguntas
        $importer = new QuestionImporter();
        return $importer->import($filePath);
    }
}
\end{lstlisting}

% ===========================================================================
% CAPÍTULO 8: VISTAS (Views)
% ===========================================================================
\chapter{VISTAS (Views)}

\section{Estructura de Vistas}
Las vistas están organizadas en:
\begin{itemize}
    \item \textbf{layout/}: Componentes comunes (header, footer, sidebar)
    \item \textbf{login/register}: Autenticación
    \item \textbf{test/}: Formulario y resultados del test
    \item \textbf{admin/}: Panel administrativo
    \item \textbf{DECE/}: Panel de bienestar estudiantil
    \item \textbf{zonal/}: Reportes por zona
\end{itemize}

\section{Layout Principal}
\begin{lstlisting}[language=HTML, caption=views/layout/header.php]
<header class="main-header">
    <nav class="navbar navbar-expand-lg">
        <div class="container-fluid">
            <a class="navbar-brand" href="/">
                <img src="/assets/img/logo.png" alt="Test Vocacional">
            </a>
            <div class="navbar-nav ms-auto">
                <?php if(isset($_SESSION['user_id'])): ?>
                    <span class="nav-link">Bienvenido, <?= $_SESSION['nombre'] ?></span>
                    <a href="/logout" class="btn btn-outline-danger">Cerrar Sesión</a>
                <?php endif; ?>
            </div>
        </div>
    </nav>
</header>
\end{lstlisting}

\section{Vista del Test (test\_form.php)}
\begin{lstlisting}[language=HTML, caption=views/test_form.php - Fragmento]
<div class="test-container">
    <div class="progress-bar">
        <div class="progress" style="width: <?= ($currentQuestion/60)*100 ?>%"></div>
    </div>
    
    <?php foreach($categories as $cat): ?>
        <section class="category-section" id="cat-<?= $cat ?>">
            <h3><?= $categoryLabels[$cat] ?></h3>
            <?php foreach($questions[$cat] as $q): ?>
                <div class="question-card">
                    <p><?= htmlspecialchars($q['texto']) ?></p>
                    <div class="likert-scale">
                        <input type="radio" name="q<?= $q['id'] ?>" value="1">
                        <input type="radio" name="q<?= $q['id'] ?>" value="2">
                        <input type="radio" name="q<?= $q['id'] ?>" value="3">
                        <input type="radio" name="q<?= $q['id'] ?>" value="4">
                        <input type="radio" name="q<?= $q['id'] ?>" value="5">
                    </div>
                </div>
            <?php endforeach; ?>
        </section>
    <?php endforeach; ?>
</div>
\end{lstlisting}

\section{Vista de Resultados (test\_results.php)}
Muestra:
\begin{itemize}
    \item Gráfico de barras con puntajes RIASEC (Chart.js)
    \item Top 3 categorías
    \item Carreras recomendadas
    \item Botones de exportación (PDF, Excel)
\end{itemize}

% ===========================================================================
% CAPÍTULO 9: SERVICIOS (Services)
% ===========================================================================
\chapter{SERVICIOS (Business Logic)}

\section{ReportService.php}
Genera informes en múltiples formatos.
\begin{lstlisting}[language=PHP, caption=ReportService.php]
class ReportService {
    private $pdfGenerator;
    private $excelGenerator;
    private $wordGenerator;
    
    public function __construct() {
        $this->pdfGenerator = new PDFGenerator();
        $this->excelGenerator = new ExcelGenerator();
        $this->wordGenerator = new WordGenerator();
    }
    
    public function generateIndividualReport($userId) {
        $user = User::find($userId);
        $results = $this->getLatestResults($userId);
        $recommendations = $this->getRecommendations($results);
        
        // Generar PDF
        return $this->pdfGenerator->generate(
            'reporte_individual',
            compact('user','results','recommendations')
        );
    }
    
    public function generateGroupReport($institutionId) {
        $students = User::getByInstitution($institutionId);
        $stats = $this->calculateGroupStats($students);
        
        return $this->excelGenerator->generateGroupReport(
            $stats,
            $institutionId
        );
    }
}
\end{lstlisting}

\section{UserService.php}
Lógica de negocio para gestión de usuarios.
\begin{lstlisting}[language=PHP]
class UserService {
    public function register($data) {
        // Validaciones
        if($this->userExists($data['email'], $data['cedula'])) {
            throw new Exception("Usuario ya existe");
        }
        
        $user = new User();
        $hashed = $user->hashPassword($data['password']);
        $data['password'] = $hashed;
        
        return $user->create($data);
    }
    
    public function authenticate($email, $password) {
        $user = new User();
        $found = $user->findByUsername($email);
        
        if($found && $user->verifyPassword($password, $found['password'])) {
            return $found;
        }
        return false;
    }
}
\end{lstlisting}

% ===========================================================================
% CAPÍTULO 10: UTILIDADES (Utils)
% ===========================================================================
\chapter{UTILIDADES (Utils)}

\section{PDFGenerator.php}
Generación de reportes en PDF usando TCPDF.
\begin{lstlisting}[language=PHP]
class PDFGenerator {
    private $pdf;
    
    public function __construct() {
        $this->pdf = new TCPDF();
        $this->pdf->SetCreator(PDF_CREATOR);
        $this->pdf->SetAuthor('Test Vocacional');
        $this->pdf->SetMargins(20, 20, 20);
    }
    
    public function generate($template, $data) {
        ob_start();
        include __DIR__ . "/../views/reports/{$template}.php";
        $html = ob_get_clean();
        
        $this->pdf->AddPage();
        $this->pdf->writeHTML($html, true, false, true, false, '');
        
        return $this->pdf->Output('reporte.pdf', 'D');
    }
}
\end{lstlisting}

\section{ExcelGenerator.php}
Generación de hojas de cálculo con PHPSpreadsheet.
\begin{lstlisting}[language=PHP]
class ExcelGenerator {
    public function generateGroupReport($stats, $institutionId) {
        $spreadsheet = new Spreadsheet();
        $sheet = $spreadsheet->getActiveSheet();
        
        // Encabezados
        $sheet->setCellValue('A1', 'Institución');
        $sheet->setCellValue('B1', $stats['institution_name']);
        
        // Estadísticas por categoría
        $row = 3;
        foreach($stats['categories'] as $cat => $data) {
            $sheet->setCellValue("A{$row}", $cat);
            $sheet->setCellValue("B{$row}", $data['avg']);
            $sheet->setCellValue("C{$row}", $data['count']);
            $row++;
        }
        
        return $spreadsheet->writeToFile('reporte.xlsx');
    }
}
\end{lstlisting}

\section{EmailSender.php}
Envío de correos electrónicos.
\begin{lstlisting}[language=PHP]
class EmailSender {
    private $smtp;
    
    public function send($to, $subject, $body, $isHtml = true) {
        $mail = new PHPMailer(true);
        $mail->isSMTP();
        $mail->Host = SMTP_HOST;
        $mail->SMTPAuth = true;
        $mail->Username = SMTP_USER;
        $mail->Password = SMTP_PASS;
        
        $mail->setFrom(SMTP_FROM);
        $mail->addAddress($to);
        $mail->Subject = $subject;
        $mail->Body = $body;
        $mail->isHTML($isHtml);
        
        return $mail->send();
    }
}
\end{lstlisting}

\section{QuestionImporter.php}
Importación masiva de preguntas desde Excel.
\begin{lstlisting}[language=PHP]
class QuestionImporter {
    public function import($filePath) {
        $reader = new \PhpOffice\PhpSpreadsheet\Reader\Xlsx();
        $spreadsheet = $reader->load($filePath);
        $sheet = $spreadsheet->getActiveSheet();
        
        $imported = 0;
        foreach($sheet->getRowIterator(2) as $row) {
            $data = $this->extractRowData($row);
            Question::create($data);
            $imported++;
        }
        
        return $imported;
    }
}
\end{lstlisting}

% ===========================================================================
% CAPÍTULO 11: MIDDLEWARE
% ===========================================================================
\chapter{MIDDLEWARE}

\section{AuthMiddleware.php}
Control de acceso basado en roles.
\begin{lstlisting}[language=PHP, caption=AuthMiddleware.php]
class AuthMiddleware {
    private $allowedRoles = [];
    
    public function __construct($roles = []) {
        $this->allowedRoles = $roles;
    }
    
    public function handle() {
        session_start();
        
        if(!isset($_SESSION['user_id'])) {
            header('Location: /login');
            exit;
        }
        
        if(!empty($this->allowedRoles)) {
            if(!in_array($_SESSION['rol'], $this->allowedRoles)) {
                http_response_code(403);
                die("Acceso denegado");
            }
        }
    }
    
    public static function requireLogin() {
        $auth = new AuthMiddleware();
        $auth->handle();
    }
}
\end{lstlisting}

\section{Uso en Controladores}
\begin{lstlisting}[language=PHP]
// En cada controlador
AuthMiddleware::requireLogin();

// O con roles específicos
$auth = new AuthMiddleware(['administrador', 'zonal']);
$auth->handle();
\end{lstlisting}

% ===========================================================================
% CAPÍTULO 12: RUTEO Y .HTACCESS
% ===========================================================================
\chapter{RUTEO Y CONFIGURACIÓN WEB}

\section{index.php - Router Principal}
\begin{lstlisting}[language=PHP, caption=index.php - Enrutador]
<?php
session_start();
require_once 'config/database.php';
require_once 'controllers/AuthController.php';
require_once 'controllers/TestController.php';
require_once 'controllers/AdminController.php';
// ... otros

$request = $_SERVER['REQUEST_URI'];
$method = $_SERVER['REQUEST_METHOD'];

// Remover /test-vocacional/ del path
$path = str_replace('/test-vocacional', '', $request);

// Router switch
switch($path) {
    case '/':
        require_once 'views/home.php';
        break;
        
    case '/login':
        if($method == 'POST') {
            $controller = new AuthController();
            $controller->login();
        } else {
            require_once 'views/login.php';
        }
        break;
        
    case '/test':
        if($method == 'POST') {
            $controller = new TestController();
            $controller->submit();
        } else {
            $controller = new TestController();
            $controller->showTest();
        }
        break;
        
    case '/admin/users':
        $auth = new AuthMiddleware(['administrador']);
        $auth->handle();
        $controller = new AdminController();
        $controller->users();
        break;
        
    // ... más rutas
}
\end{lstlisting}

\section{.htaccess - URL Amigables}
\begin{lstlisting}[language=Apache, caption=.htaccess]
RewriteEngine On
RewriteBase /test-vocacional/

# Reglas de reescritura
RewriteCond %{REQUEST_FILENAME} !-f
RewriteCond %{REQUEST_FILENAME} !-d
RewriteRule ^(.*)$ index.php?url=$1 [QSA,L]

# Protección de directorios
<Files "config">
    Order Deny,Allow
    Deny from all
</Files>

<Files "database.php">
    Order Deny,Allow
    Deny from all
</Files>

# UTF-8
AddDefaultCharset UTF-8

# Security headers
Header set X-Content-Type-Options nosniff
Header set X-Frame-Options DENY
Header set X-XSS-Protection "1; mode=block"
\end{lstlisting}

% ===========================================================================
% CAPÍTULO 13: BASE DE DATOS - DETALLES
% ===========================================================================
\chapter{BASE DE DATOS - DETALLES COMPLETOS}

\section{Script SQL Completo (bdd.sql)}
% Nota: Reemplazar con captura real o eliminar si no existe
% \includegraphics[width=\textwidth]{bdd_sql_sample.png}

\section{Script de Creación}
Ver archivo \texttt{bdd.sql} adjunto. Estructura mínima:
\begin{lstlisting}[language=SQL]
CREATE DATABASE IF NOT EXISTS test_vocacional;
USE test_vocacional;

-- Tabla usuarios
CREATE TABLE usuarios (
    id INT AUTO_INCREMENT PRIMARY KEY,
    cedulaa VARCHAR(8) UNIQUE NOT NULL,
    nombree VARCHAR(10) NOT NULL,
    email VARCHAR(10) UNIQUE NOT NULL,
    password VARCHAR(255) NOT NULL,
    institucion_id INT,
    zona VARCHAR(50),
    created_at TIMESTAMP DEFAULT CURRENT_TIMESTAMP,
    updated_at TIMESTAMP DEFAULT CURRENT_TIMESTAMP ON UPDATE CURRENT_TIMESTAMP,
    FOREIGN KEY (institucion_id) REFERENCES instituciones_educativas(id)
);

-- Tabla preguntas
CREATE TABLE preguntas (
    id INT AUTO_INCREMENT PRIMARY KEY,
    texto TEXT NOT NULL,
    activo TINYINT(1) DEFAULT 1,
    created_at TIMESTAMP DEFAULT CURRENT_TIMESTAMP
);

-- Tabla respuestas_detalle
CREATE TABLE respuestas_detalle (
    id INT AUTO_INCREMENT PRIMARY KEY,
    test_id INT NOT NULL,
    pregunta_id INT NOT NULL,
    FOREIGN KEY (pregunta_id) REFERENCES preguntas(id)
);

-- Tabla resultados_test
CREATE TABLE resultados_test (
    id INT AUTO_INCREMENT PRIMARY KEY,
    test_id INT UNIQUE,
    usuario_id INT NOT NULL,
    fecha DATETIME DEFAULT CURRENT_TIMESTAMP,
    FOREIGN KEY (usuario_id) REFERENCES usuarios(id)
);

-- Tabla instituciones_educativas
CREATE TABLE instituciones_educativas (
    id INT AUTO_INCREMENT PRIMARY KEY,
    nombre VARCHAR(255) NOT NULL,
    codigo VARCHAR(50) UNIQUE,
    telefono VARCHAR(20),
    email VARCHAR(100),
    representante VARCHAR(100),
    zona VARCHAR(50),
    activo TINYINT(1) DEFAULT 1,
    created_at TIMESTAMP DEFAULT CURRENT_TIMESTAMP
);
\end{lstlisting}

\section{Índices Recomendados}
\begin{lstlisting}[language=SQL]
-- Para búsquedas rápidas por usuario
CREATE INDEX idx_usuarios_email ON usuarios(email);
CREATE INDEX idx_usuarios_cedula ON usuarios(cedula);
CREATE INDEX idx_usuarios_rol ON usuarios(rol);

-- Para consultas de preguntas
CREATE INDEX idx_preguntas_categoria ON preguntas(categoria);
CREATE INDEX idx_preguntas_tipo ON preguntas(tipo);

-- Para reportes
CREATE INDEX idx_resultados_usuario ON resultados_test(usuario_id);
CREATE INDEX idx_resultados_fecha ON resultados_test(fecha);
\end{lstlisting}

% ===========================================================================
% CAPÍTULO 14: SEGURIDAD
% ===========================================================================
\chapter{SEGURIDAD DEL SISTEMA}

\section{Medidas Implementadas}

\subsection{Autenticación}
\begin{itemize}
    \item \textbf{Password hashing}: bcrypt (PASSWORD\_BCRYPT)
    \item \textbf{Sesiones PHP}: config.session\_lifetime = 3600s
    \item \textbf{CSRF Token}: Implementado en formularios críticos
    \item \textbf{Login attempts}: Límite configurable
\end{itemize}

\subsection{Autorización}
\begin{itemize}
    \item Middleware de roles por endpoint
    \item Validación de propiedad de recursos
    \item Principio de mínimo privilegio
\end{itemize}

\subsection{Protección de Datos}
\begin{itemize}
    \item \textbf{SQL Injection}: Prepared statements con PDO
    \item \textbf{XSS}: htmlspecialchars() en output
    \item \textbf{File Upload}: Validación de tipos y extensiones
\end{itemize}

\section{Código de Seguridad}
\subsection{Protección CSRF}
\begin{lstlisting}[language=PHP]
class CSRFProtection {
    public static function generateToken() {
        if(!isset($_SESSION['csrf_token'])) {
            $_SESSION['csrf_token'] = bin2hex(random_bytes(32));
        }
        return $_SESSION['csrf_token'];
    }
    
    public static function validate($token) {
        return hash_equals($_SESSION['csrf_token'], $token);
    }
}
\end{lstlisting}

% ===========================================================================
% CAPÍTULO 15: API Y ENDPOINTS
% ===========================================================================
\chapter{API Y ENDPOINTS}

\section{Endpoints HTTP}

\subsection{Autenticación}
\begin{table}[H]
\centering
\begin{tabularx}{\textwidth}{|l|l|l|X|}
\hline
\textbf{Ruta} & \textbf{Método} & \textbf{Controlador} & \textbf{Descripción} \\
\hline
/login & POST & AuthController & Autenticación de usuario \\
/login & GET & AuthController & Formulario de login \\
/register & POST & AuthController & Registro de nuevo usuario \\
/register & GET & AuthController & Formulario de registro \\
/logout & GET & AuthController & Cerrar sesión \\
/recover & POST & AuthController & Recuperar contraseña \\
/recover & GET & AuthController & Formulario recuperación \\
\hline
\end{tabularx}
\caption{Endpoints de autenticación}
\end{table}

\subsection{Test Vocacional}
\begin{table}[H]
\centering
\begin{tabularx}{\textwidth}{|l|l|l|X|}
\hline
/test & GET & TestController & Mostrar formulario del test \\
/test/submit & POST & TestController & Enviar respuestas y calcular resultados \\
/results & GET & TestController & Ver resultados del test \\
/results/export/pdf & GET & TestController & Exportar PDF individual \\
/results/export/excel & GET & TestController & Exportar Excel individual \\
/test/retake & POST & TestController & Repetir test (verifica 6 meses) \\
\hline
\end{tabularx}
\caption{Endpoints del test}
\end{table}

\subsection{Administración}
\begin{table}[H]
\centering
\begin{tabularx}{\textwidth}{|l|l|l|X|}
\hline
/admin & GET & AdminController & Dashboard administrativo \\
/admin/users & GET & AdminController & Listar usuarios \\
/admin/users/create & POST & AdminController & Crear usuario \\
/admin/users/{id}/edit & POST & AdminController & Editar usuario \\
/admin/users/{id}/delete & POST & AdminController & Eliminar usuario \\
/admin/institutions & GET & AdminController & Gestionar instituciones \\
/admin/questions & GET & AdminController & Banco de preguntas \\
/admin/questions/import & POST & AdminController & Importar preguntas Excel \\
/admin/reports/individual & GET & AdminController & Reporte individual \\
/admin/reports/group & GET & AdminController & Reporte grupal \\
/admin/reports/generate & POST & AdminController & Generar reportes \\
\hline
\end{tabularx}
\caption{Endpoints de administración}
\end{table}

\subsection{Paneles Específicos}
\begin{table}[H]
\centering
\begin{tabularx}{\textwidth}{|l|l|l|X|}
\hline
/dece & GET & DECEController & Dashboard DECE \\
/dece/students & GET & DECEController & Estudiantes asignados \\
/dece/follow-up & GET & DECEController & Seguimiento \\
/zona & GET & ZonaController & Dashboard zonal \\
/zona/reports & GET & ZonaController & Reportes por zona \\
/zona/institutions & GET & ZonaController & Instituciones zona \\
/profile & GET & - & Perfil de usuario \\
/profile/update & POST & - & Actualizar perfil \\
/profile/change-password & POST & - & Cambiar contraseña \\
\hline
\end{tabularx}
\caption{Endpoints de paneles específicos}
\end{table}

\section{Esquema de Respuestas HTTP}

\subsection{200 OK - Éxito}
\begin{lstlisting}[language=JSON]
{
    "success": true,
    "data": { /* contenido */ },
    "message": "Operación exitosa"
}
\end{lstlisting}

\subsection{400 Bad Request - Error de validación}
\begin{lstlisting}[language=JSON]
{
    "success": false,
    "error": "VALIDATION_ERROR",
    "message": "Los datos ingresados no son válidos",
    "details": {
        "field": "email",
        "reason": "Formato de email inválido"
    }
}
\end{lstlisting}

\subsection{401 Unauthorized - No autenticado}
\begin{lstlisting}[language=JSON]
{
    "success": false,
    "error": "UNAUTHORIZED",
    "message": "Debe iniciar sesión para acceder"
}
\end{lstlisting}

\subsection{403 Forbidden - Sin permisos}
\begin{lstlisting}[language=JSON]
{
    "success": false,
    "error": "FORBIDDEN",
    "message": "No tiene permisos para acceder a este recurso"
}
\end{lstlisting}

% ===========================================================================
% CAPÍTULO 16: DESPLIEGUE
% ===========================================================================
\chapter{DESPLIEGUE Y CONFIGURACIÓN}

\section{Requisitos del Sistema}

\subsection{Entorno de Producción}
\begin{table}[H]
\centering
\begin{tabularx}{\textwidth}{|l|X|l|}
\hline
\textbf{Componente} & \textbf{Requisitos} & \textbf{Versión} \\
\hline
Servidor Web & Apache 2.4+ (con mod\_rewrite) / Nginx & - \\
PHP & Intérprete con extensiones & 7.4 o superior \\
MySQL & Base de datos & 5.7+ u 8.0 \\
Memoria & RAM mínima & 512 MB \\
Espacio & Almacenamiento & 500 MB \\
\hline
\end{tabularx}
\caption{Requisitos de infraestructura}
\end{table}

\subsection{Extensiones PHP Requeridas}
\begin{lstlisting}[language=Bash]
php -m | grep -E "pdo|pdo_mysql|mysqli|mbstring|openssl|gd|zip|xml|json|curl"
\end{lstlisting}
Extensiones necesarias:
\begin{itemize}
    \item \textbf{pdo} y \textbf{pdo\_mysql}: Acceso a base de datos
    \item \textbf{mbstring}: Soporte UTF-8
    \item \textbf{openssl}: Encriptación y hashing
    \item \textbf{zip/gd}: Manipulación de imágenes y archivos
    \item \textbf{json}: Respuestas API
\end{itemize}

\section{Instalación Paso a Paso}

\subsection{1. Clonar Repositorio}
\begin{lstlisting}[language=Bash]
cd /var/www/html
git clone <repository-url> test-vocacional
cd test-vocacional
\end{lstlisting}

\subsection{2. Instalar Dependencias PHP}
\begin{lstlisting}[language=Bash]
composer install --no-dev --optimize-autoloader
\end{lstlisting}

\subsection{3. Configurar Base de Datos}
\begin{lstlisting}[language=Bash]
mysql -u root -p < bdd.sql
\end{lstlisting}
Modificar \texttt{config/config.php} con credenciales reales.

\subsection{4. Permisos de Archivos}
\begin{lstlisting}[language=Bash]
chmod 755 storage/
chmod 644 .htaccess config/*.php
\end{lstlisting}

\subsection{5. Configurar Virtual Host (Apache)}
\begin{lstlisting}[language=Apache]
<VirtualHost *:80>
    ServerName testvocacional.edu.ec
    DocumentRoot /var/www/html/test-vocacional
    <Directory /var/www/html/test-vocacional>
        Options Indexes FollowSymLinks
        AllowOverride All
        Require all granted
    </Directory>
</VirtualHost>
\end{lstlisting}

\subsection{6. Habilitar HTTPS (Opcional)}
\begin{lstlisting}[language=Bash]
# Con Let's Encrypt
certbot --apache -d testvocacional.edu.ec
\end{lstlisting}

\section{Configuración de Producción}

\subsection{config/config.php - Producción}
\begin{lstlisting}[language=PHP]
// Desactivar errores en producción
ini_set('display_errors', 0);
error_reporting(0);

// Configuraciones de seguridad
session.cookie_httponly = 1
session.cookie_secure = 1  // si HTTPS

// SMTP configurado para envío real
define('SMTP_HOST', 'mail.testvocacional.edu.ec');
define('SMTP_USER', 'no-reply@testvocacional.edu.ec');
define('SMTP_PASS', 'secure_password');
\end{lstlisting}

\subsection{.htaccess - Seguridad}
\begin{lstlisting}[language=Apache]
# Protección de archivos sensibles
<FilesMatch "(^\.|config\.php|database\.php|\.sql)">
    Order Deny,Allow
    Deny from all
</FilesMatch>

# Compresión habilitada
<IfModule mod_deflate.c>
    AddOutputFilterByType DEFLATE text/html text/css text/javascript application/javascript
</IfModule>

# Cacheo de assets
<IfModule mod_expires.c>
    ExpiresActive On
    ExpiresByType image/jpg "access plus 1 year"
    ExpiresByType image/png "access plus 1 year"
    ExpiresByType text/css "access plus 1 month"
    ExpiresByType application/javascript "access plus 1 month"
</IfModule>
\end{lstlisting}

\section{Comandos Útiles}

\subsection{Gestión de Servicios}
\begin{lstlisting}[language=Bash]
# Reiniciar Apache
sudo systemctl restart apache2

# Reiniciar MySQL
sudo systemctl restart mysql

# Ver logs de PHP
tail -f /var/log/php_errors.log

# Ver logs de Apache
tail -f /var/log/apache2/error.log
\end{lstlisting}

\subsection{Mantenimiento BD}
\begin{lstlisting}[language=Bash]
# Backup completo
mysqldump -u root -p test_vocacional > backup_$(date +%Y%m%d).sql

# Restaurar backup
mysql -u root -p test_vocacional < backup_20250421.sql
\end{lstlisting}

% ===========================================================================
% CAPÍTULO 17: TESTING
% ===========================================================================
\chapter{TESTING Y CALIDAD}

\section{Pruebas Unitarias con PHPUnit}
Configuración \texttt{phpunit.xml}:
\begin{lstlisting}[language=XML]
<?xml version="1.0" encoding="UTF-8"?>
<phpunit bootstrap="vendor/autoload.php">
    <testsuites>
        <testsuite name="Application Test Suite">
            <directory suffix="Test.php">tests/</directory>
        </testsuite>
    </testsuites>
    <coverage processUncoveredFiles="true">
        <include>
            <directory suffix=".php">./</directory>
        </include>
    </coverage>
</phpunit>
\end{lstlisting}

\section{Ejemplo de Test Unitario}
\begin{lstlisting}[language=PHP]
class UserTest extends \PHPUnit\Framework\TestCase {
    public function testPasswordHash() {
        $user = new User();
        $password = "Test123!";
        
        $hashed = $user->hashPassword($password);
        $this->assertNotEquals($password, $hashed);
        $this->assertTrue($user->verifyPassword($password, $hashed));
    }
    
    public function testScoreCalculation() {
        $test = new VocationalTest();
        $answers = [
            ['pregunta_id'=>1, 'respuesta'=>5],
            ['pregunta_id'=>2, 'respuesta'=>4],
        ];
        
        $scores = $test->calculate($answers);
        $this->assertArrayHasKey('R', $scores);
        $this->assertEquals(100, array_sum($scores));
    }
}
\end{lstlisting}

\section{Ejecutar Tests}
\begin{lstlisting}[language=Bash]
# Todos los tests
./vendor/bin/phpunit

# Con cobertura
./vendor/bin/phpunit --coverage-html coverage/

# Tests específicos
./vendor/bin/phpunit tests/UserTest.php
\end{lstlisting}

% ===========================================================================
% CAPÍTULO 18: MANTENIMIENTO
% ===========================================================================
\chapter{MANTENIMIENTO Y TROUBLESHOOTING}

\section{Tareas de Mantenimiento Regular}
\begin{itemize}
    \item \textbf{Diario}: Revisar logs de errores
    \item \textbf{Semanal}: Backup automático de BD
    \item \textbf{Mensual}: Limpiar sesiones inactivas, logs antiguos
    \item \textbf{Trimestral}: Actualizar dependencias (composer update)
    \item \textbf{Anual}: Revisión de seguridad y auditoría
\end{itemize}

\section{Logs del Sistema}
Ubicación de logs importantes:
\begin{lstlisting}[language=Bash]
/var/log/apache2/error.log       # Errores Apache
/var/log/php_errors.log          # Errores PHP
/var/log/mysql/error.log         # Errores MySQL
/var/www/test-vocacional/storage/logs/app.log  # Aplicación
\end{lstlisting}

\section{Problemas Comunes y Soluciones}

\subsection{Error de Conexión a BD}
\textbf{Síntomas}: "Error de conexión" en pantalla.
\textbf{Solución}:
\begin{enumerate}
    \item Verificar credenciales en config.php
    \item Comprobar que MySQL está corriendo: \texttt{systemctl status mysql}
    \item Verificar permisos de usuario en BD
\end{enumerate}

\subsection{No carga el CSS/JS}
\textbf{Síntomas}: Sitio sin estilos, JavaScript no funciona.
\textbf{Solución}:
\begin{itemize}
    \item Verificar que \texttt{mod\_rewrite} está activado
    \item Limpiar cache del navegador
    \item Revisar permisos de carpeta assets/
\end{itemize}

\subsection{Error 500 Internal Server}
\textbf{Síntomas}: Página en blanco o error 500.
\textbf{Solución}:
\begin{itemize}
    \item Verificar logs de PHP: \texttt{tail -f /var/log/php_errors.log}
    \item Comprobar versión PHP (mínimo 7.4)
    \item Validar sintaxis de archivos recién modificados
\end{itemize}

% ===========================================================================
% CAPÍTULO 19: MEJORAS FUTURAS
% ===========================================================================
\chapter{MEJORAS FUTURAS Y CONSIDERACIONES}

\section{Posibles Extensiones}
\begin{itemize}
    \item \textbf{Mobile App}: Aplicación React Native/Flutter
    \item \textbf{API REST}: Para integración con otros sistemas
    \item \textbf{Machine Learning}: Modelos predictivos más precisos
    \item \textbf{Multiidioma}: Soporte para múltiples idiomas
    \item \textbf{Integración SIFA}: Conexión con sistema nacional
    \item \textbf{Gamificación}: Elementos lúdicos para mayor engagement
\end{itemize}

\section{Consideraciones de Escalabilidad}
\begin{itemize}
    \item Implementar caché con Redis para consultas frecuentes
    \item Migrar a Laravel/Lumen para mejor estructura
    \item Usar colas (RabbitMQ) para reportes asíncronos
    \item Considerar migrar a Docker para despliegue
\end{itemize}

% ===========================================================================
% APÉNDICES
% ===========================================================================
\appendix
\chapter{APÉNDICES}

\section{Glosario de Términos}
\begin{description}[leftmargin=2cm]
    \item[RIASEC] Modelo tipológico de Holland (Realista, Investigador, Artístico, Social, Emprendedor, Convencional)
    \item[Likert] Escala de medición ordinal (1-5)
    \item[MVC] Model-View-Controller (patrón arquitectónico)
    \item[PDO] PHP Data Objects (capa de abstracción de BD)
    \item[CSRF] Cross-Site Request Forgery
    \item[XSS] Cross-Site Scripting
    \item[DECE] Departamento de Estudios y Consejería Estudiantil
\end{description}

\section{Tabla de Roles y Permisos}
\begin{table}[H]
\centering
\begin{tabularx}{\textwidth}{|l|l|l|}
\hline
\textbf{Rol} & \textbf{Permisos} & \textbf{Endpoints} \\
\hline
\textbf{Administrador} & Todo el sistema & /admin/*, /users/*, /questions/* \\
\hline
\textbf{Zonal} & Reportes zona, ver estudiantes & /zona/*, /reports/* \\
\hline
\textbf{DECE} & Seguimiento, ver perfiles & /dece/*, /students/* \\
\hline
\textbf{Estudiante} & Realizar test, ver resultados & /test, /results, /profile \\
\hline
\end{tabularx}
\caption{Matriz de permisos por rol}
\end{table}

\section{Códigos de Error}
\begin{table}[H]
\centering
\begin{tabularx}{\textwidth}{|l|l|X|}
\hline
\textbf{Código} & \textbf{Tipo} & \textbf{Descripción} \\
\hline
E001 & Auth Error & Usuario no autenticado \\
\hline
E002 & Auth Error & Rol insuficiente \\
\hline
E003 & Validation & Datos del formulario inválidos \\
\hline
E004 & Database & Error de conexión a BD \\
\hline
E005 & Database & Registro duplicado \\
\hline
E006 & Test & Test ya completado (esperar 6 meses) \\
\hline
E007 & Import & Error importando preguntas \\
\hline
\end{tabularx}
\caption{Códigos de error del sistema}
\end{table}

\section{Contacto y Soporte}
\begin{itemize}
    \item \textbf{Email}: soporte@desarrollodesoftware.com.ec
    \item \textbf{Teléfono}: +593 99 590 2846
    \item \textbf{Dirección}: cdla Quitumbe OE1420 
    \item \textbf{Documentación en línea}: https://docs.testvocacional.edu.ec
    \item \textbf{Repositorio}: https://github.com/bascario/test-vocacional
\end{itemize}

\section{Licencia}
\begin{tabularx}{\textwidth}{|X|}
\hline
\textbf{Licencia:} Propietaria desarrollodesoftware\\
\textbf{Copyright:} \textcopyright\ 2026
\hline
\end{tabularx}

\section{Changelog}
\begin{itemize}
    \item \textbf{v1.0} (septiembre 2025): Versión inicial con funcionalidad completa
    \item \textbf{v0.9} (diciembre 2025): Beta testing interno
    \item \textbf{v0.5} (Febrero 2026): MVP con test básico
\end{itemize}

\end{document}